\documentclass[11pt]{article}
\usepackage{estudio}
\cuadernillo{1 --- Introducción}

\begin{document}

\portada{1}{Introducción}{%
  Diapositivas 1 a 4 \textbullet\ expositor A \textbullet\ 1:45\\[4pt]
  Sistema real, materia activa y modelo matemático}

\tableofcontents
\clearpage

% ===========================================================================
\section{Cómo se usa esta guía}

Son cuatro cuadernillos, uno por sección de la presentación. Cada uno tiene la
misma estructura:

\begin{enumerate}
  \item \textbf{Guion}: qué decir en cada diapositiva y cuánto tiempo tenés.
  \item \textbf{El fondo}: la teoría que sostiene lo que estás diciendo.
  \item \textbf{Preguntas y respuestas}: lo que te pueden preguntar y qué contestar.
\end{enumerate}

El reparto acordado es: \textbf{A} expone Introducción y después los resultados
de Vicsek; \textbf{B} expone Implementación y después los del votante;
\textbf{C} expone Simulaciones y después las comparaciones, el benchmark y las
conclusiones. La rotación es A\,$\to$\,B\,$\to$\,C\,$\to$\,A\,$\to$\,B\,$\to$\,C.

\ojo{La guía de la cátedra (punto 3.1) dice que \emph{todos} los miembros del
grupo deben estar en condiciones de exponer la presentación completa, y que si
hay preguntas al final las responden en orden, sin superponerse (punto 3.2).
Por eso los cuatro cuadernillos los leen los tres, aunque cada uno exponga solo
su parte.}

\subsection{El presupuesto de tiempo}

La presentación tiene 28 páginas: 1 de título, 5 divisorias de sección,
1 de cierre y \textbf{21 de contenido}. Trece minutos son 780 segundos, así que
el promedio real es de unos \textbf{36 segundos por diapositiva de contenido}.

\dato{Ese número es el que manda. Una diapositiva de gráfico bien explicada
lleva 30--35 s; una de animación, 40--45 s; una de ecuaciones, 50--60 s. No hay
margen para introducir, contextualizar ni repetir. Si te pasás 10 s en cada una
de las primeras diez diapositivas, llegás a conclusiones con dos minutos menos
de los que necesitás.}

Distribución cerrada, ya descontado el tiempo de las divisorias:

\begin{center}
\begin{tabular}{llcc}
\toprule
Expositor & Bloque & Diapositivas & Tiempo \\
\midrule
A & Introducción            & 1--4   & 1:45 \\
B & Implementación          & 5--7   & 1:37 \\
C & Simulaciones            & 8--10  & 1:35 \\
A & Resultados: Vicsek      & 11--16 & 2:40 \\
B & Resultados: votante     & 17--21 & 2:30 \\
C & Comparaciones y cierre  & 22--28 & 2:40 \\
\midrule
\multicolumn{3}{l}{\textbf{Total}} & \textbf{12:47} \\
\bottomrule
\end{tabular}
\end{center}

Quedan 13 segundos de colchón sobre los 13:00. Es poquísimo: ensayen con
cronómetro y con las transiciones entre personas incluidas, que es donde se
pierden los segundos.

% ===========================================================================
\section{Guion: diapositivas 1 a 4}

\diapo{1}{Portada}{15 s}{A}

Nombre del trabajo, los tres integrantes, materia y comisión. No leas la
diapositiva entera.

\begin{quote}
\itshape
``Buenos días. Trabajo Práctico 2 de Simulación de Sistemas, grupo 5. Vamos a
presentar un autómata celular off-lattice de bandadas de agentes
autopropulsados, comparando el modelo estándar de Vicsek contra el modelo de
votante. Arranco yo con la introducción.''
\end{quote}

\diapo{3}{Sistema real}{38 s}{A}

La foto es una murmuración de estorninos. Los cuatro \textit{bullets} de la
diapositiva son el orden en que hay que decirlo, y cada uno es una frase, no un
párrafo.

\begin{quote}
\itshape
``El sistema real son las bandadas: estorninos, cardúmenes de peces, colonias
bacterianas. Lo característico es que \emph{no hay un líder}: cada individuo
decide su dirección mirando únicamente a los vecinos que tiene al lado, y del
comportamiento local emerge un orden global coherente. Esto es un ejemplo de
materia activa, que es todo sistema donde la energía entra localmente, en cada
agente, y no desde un forzado externo.''
\end{quote}

\clave{Los tres conceptos que tenés que dejar plantados acá, porque los usan las
otras tres secciones: \textbf{(1)} interacción puramente local,
\textbf{(2)} orden global emergente sin coordinación central,
\textbf{(3)} materia activa = energía inyectada en cada partícula.}

\ojo{No te extiendas con la biología. La cátedra pide una \emph{somera}
descripción del sistema real (guía 2.1, máximo 3 diapositivas). El contenido
técnico va en las secciones siguientes; acá solo estás situando el problema.}

\diapo{4}{Modelo matemático}{50 s}{A}

Es la diapositiva más densa de tu bloque: tres ecuaciones y la definición del
ruido. Ordená así:

\begin{enumerate}
  \item \textbf{Vicsek}: cada partícula adopta la dirección \emph{promedio} de
    sus vecinos dentro del radio de interacción, más un ruido angular.
  \item \textbf{Votante}: no promedia; elige \emph{un} vecino al azar y copia su
    dirección, más el mismo ruido.
  \item \textbf{Posición}: común a los dos modelos, y la rapidez es constante
    ---lo único que cambia es la dirección.
  \item \textbf{Ruido}: uniforme en $[-\eta/2,\ \eta/2]$. $\eta$ es el único
    parámetro de control.
  \item \textbf{Adimensionalidad}: $\Delta t$, $v$ y $\eta$ no tienen unidades
    físicas.
\end{enumerate}

\begin{quote}
\itshape
``El modelo matemático de Vicsek, de 1995, dice que cada partícula toma la
dirección promedio de las que tiene dentro de un radio de interacción, más un
ruido angular. El promedio se escribe como el arcotangente del cociente entre el
seno medio y el coseno medio. El modelo de votante cambia una sola cosa: en vez
de promediar, la partícula elige un vecino al azar con distribución uniforme y
copia su dirección, con el mismo ruido. La posición se actualiza igual en los
dos casos, con rapidez constante: lo único que evoluciona es la dirección. El
ruido es uniforme entre menos eta medios y eta medios, y es el único parámetro
de control del desorden. Todas las magnitudes son adimensionales.''
\end{quote}

\clave{La frase que tiene que quedar sonando, porque es la tesis de todo el
trabajo: \emph{``la única diferencia entre los dos modelos es promediar sobre
todos los vecinos versus copiar a uno solo elegido al azar''}. Todo lo que
mostramos después se explica con eso.}

\ojo{Acá \emph{no} se habla de $L=10$, ni de $\rho=2,4,8$, ni del Cell Index
Method. La guía 2.1 es explícita: en Introducción va el modelo matemático
general, ``sin especificar ningún sistema particular, ni métodos sobre cómo se
resuelve el modelo''. Los valores concretos son de C (diapositiva 9) y el método
es de B (diapositiva 7). Si te adelantás, le sacás contenido a tus compañeros y
la cátedra lo penaliza.}

Cerrá pasándole la palabra a B:

\begin{quote}
\itshape
``Con el modelo matemático planteado, [B] sigue con cómo lo llevamos a un
modelo computacional.''
\end{quote}

% ===========================================================================
\section{El fondo: lo que hay detrás de esas cuatro diapositivas}

\subsection{Por qué esto es un autómata celular, y por qué \textit{off-lattice}}

El TP se titula ``Autómatas Celulares''. Un autómata celular clásico tiene tres
ingredientes: un conjunto de agentes con estado discreto, una regla de
actualización \emph{local} (el estado nuevo depende solo de un vecindario), y
actualización \emph{sincrónica} (todos cambian a la vez, mirando la misma foto
del sistema).

Nuestro modelo tiene los tres, pero con una diferencia: las partículas no viven
en las celdas fijas de una grilla, sino en posiciones continuas del plano, y se
\emph{mueven}. Por eso es \textit{off-lattice}. La consecuencia práctica es
importante: el vecindario de una partícula no es fijo, se redefine en cada paso
según quién esté a distancia menor que $r_c$. Esa es exactamente la razón por la
que hace falta el Cell Index Method y por la que la grilla se reconstruye en
cada paso (diapositiva 7).

\teoria{En un autómata celular de red, el vecindario de una celda es siempre el
mismo (los 4 u 8 de al lado). Acá el vecindario es dinámico y hay que
recalcularlo, y eso es lo que convierte la búsqueda de vecinos en el cuello de
botella del motor. El TP1 existió justamente para resolver ese problema.}

\subsection{Materia activa}

Un fluido común se mueve porque algo lo empuja desde afuera. En materia activa
cada partícula tiene su propio motor: consume energía y la convierte en
movimiento dirigido. Eso rompe el balance detallado y saca al sistema del
equilibrio termodinámico, lo que permite fenómenos que en equilibrio están
prohibidos ---como el orden orientacional de largo alcance en dos dimensiones,
que es precisamente lo que observa Vicsek.

En nuestro modelo la actividad está en que $|\vect{v}_i| = v$ es constante y no
nula: la partícula nunca se frena, siempre está inyectando movimiento. El ruido
$\eta$ hace el papel de la temperatura, y la regla de alineación hace el papel
de la interacción atractiva entre espines.

\subsection{La analogía con el modelo XY, y por qué importa}

La regla de Vicsek es, formalmente, la de un ferromagneto XY: cada ``espín''
(la dirección $\theta_i$) tiende a alinearse con sus vecinos, y la temperatura
(el ruido $\eta$) lo desordena. La diferencia crucial es que los espines de un
XY están clavados en una red, y los nuestros \emph{se mueven}.

Eso no es un detalle. El teorema de Mermin--Wagner prohíbe el orden de largo
alcance en un XY bidimensional a temperatura finita. Vicsek encontró que su
modelo \emph{sí} ordena en 2D, y la razón es el movimiento: las partículas
transportan información de orientación a través del sistema, generando una
interacción de alcance efectivo mucho mayor que $r_c$.

\teoria{Esta es la razón por la que el paper de Vicsek de 1995 se llama
``\textit{Novel type of phase transition in a system of self-driven
particles}''. Lo novedoso no es que haya transición de fase, es que la haya
\emph{en dos dimensiones}, donde el equilibrio la prohíbe.}

\subsection{Transición de fase y parámetro de orden}

Un parámetro de orden es una cantidad que vale (aproximadamente) cero en la fase
desordenada y algo distinto de cero en la ordenada. Acá es la polarización
$v_a$, que define C en la diapositiva 10. Al aumentar $\eta$ el sistema pasa de
una bandada coherente ($v_a \approx 1$) a movimiento incoherente
($v_a \approx 0$): eso es la transición orden--desorden que caracterizamos en
todo el trabajo.

\subsection{Por qué el ruido es uniforme y no gaussiano}

Es la convención del paper original de Vicsek: $\Delta\theta$ uniforme en
$[-\eta/2,\ \eta/2]$. Tiene dos ventajas prácticas. Primero, $\eta$ queda
acotado: en $\eta = 2\pi$ el ruido cubre toda la circunferencia y el sistema
está máximamente desordenado, lo que da un rango de barrido natural y cerrado,
$\eta \in [0,\ 2\pi]$. Segundo, es directamente comparable con la bibliografía.
Con un ruido gaussiano habría que elegir un corte arbitrario y el rango no
tendría un extremo natural.

\subsection{Por qué el modelo de votante}

Viene de Loscar, Baglietto y Vázquez (2021). El modelo de votante clásico es un
modelo de opinión: un individuo adopta la opinión de un vecino al azar. Acá la
``opinión'' es la dirección de movimiento, y es \emph{multiestado} porque
$\theta$ es continuo en vez de binario.

La diferencia con Vicsek es de mecanismo, no de grado, y esto conviene tenerlo
claro desde la introducción porque es lo que explica todos los resultados:

\begin{itemize}
  \item \textbf{Vicsek promedia}: el promedio de $k$ direcciones ruidosas
    atenúa el ruido individual por un factor $\sim k^{-1/2}$. Más vecinos
    $\Rightarrow$ más robusto al ruido.
  \item \textbf{El votante copia}: transmite el ruido íntegro, sin atenuar. La
    cantidad de vecinos no ayuda a filtrar nada.
\end{itemize}

\clave{Si te preguntan en la introducción ``¿qué esperan que pase?'', la
respuesta correcta ---y que después los datos confirman--- es: \emph{el votante
debería perder el orden a ruidos mucho más chicos que Vicsek, porque no atenúa
el ruido}. Y efectivamente: Vicsek cruza $v_a = 0{,}5$ en $\eta \approx 2{,}9$ a
$\rho=2$, y el votante en $\eta \approx 0{,}37$. Casi un orden de magnitud.}

% ===========================================================================
\section{Preguntas y respuestas}

\subsection{Sobre el sistema y el modelo}

\pregunta{¿Por qué esto es un autómata celular si las partículas no están en una
red?}
\respuesta{Porque conserva las tres propiedades que definen un autómata celular:
estado por agente, regla de actualización local sobre un vecindario, y
actualización sincrónica. Lo que cambia es que el vecindario no es fijo: las
partículas viven en posiciones continuas y se mueven, así que el vecindario se
recalcula en cada paso según quién esté a distancia menor que $r_c$. Por eso se
lo llama autómata off-lattice.}

\pregunta{¿Qué es materia activa y por qué este sistema lo es?}
\respuesta{Es un sistema donde la energía se inyecta localmente, en cada agente,
en lugar de venir de un forzado externo. Acá cada partícula se mueve con rapidez
constante y no nula: tiene su propio motor. Eso saca al sistema del equilibrio
termodinámico, y es lo que permite que haya orden de largo alcance en dos
dimensiones.}

\pregunta{¿Por qué el ruido es uniforme y no gaussiano?}
\respuesta{Es la convención del paper original de Vicsek, y hace que $\eta$
tenga un rango natural cerrado: en $\eta = 2\pi$ el ruido cubre toda la
circunferencia y el sistema está máximamente desordenado. Por eso barrimos
$\eta$ entre $0$ y $2\pi$. Con un ruido gaussiano habría que elegir un corte
arbitrario.}

\pregunta{¿Por qué el ruido va entre $-\eta/2$ y $+\eta/2$ y no entre $-\eta$ y
$+\eta$?}
\respuesta{Para que $\eta$ sea directamente la \emph{amplitud} total del ruido:
el ancho del intervalo es exactamente $\eta$. Así, $\eta = 2\pi$ significa que
el ángulo nuevo es uniforme sobre toda la circunferencia, que es el desorden
máximo. Es la convención de Vicsek 1995.}

\pregunta{¿Por qué la rapidez es constante? ¿No sería más realista que varíe?}
\respuesta{Es una simplificación deliberada del modelo mínimo: se quiere aislar
el efecto del alineamiento direccional, así que se fija el módulo de la
velocidad y se deja evolucionar solo la dirección. Si la rapidez también
variara, habría un segundo mecanismo de ordenamiento mezclado y no se podría
atribuir la transición únicamente a la competencia entre alineación y ruido.}

\pregunta{¿Qué unidades tienen $\eta$, $v$ y $\Delta t$?}
\respuesta{Ninguna, son adimensionales. El modelo no está calibrado contra un
sistema físico concreto: es un modelo mínimo. Lo único que importa son las
razones entre escalas ---por ejemplo, que $v\,\Delta t = 0{,}03$ sea el 3\% de
$r_c$, o sea que una partícula se mueva poco respecto de su radio de
interacción en un paso.}

\pregunta{¿Por qué $\eta$ llega hasta $2\pi$ y no más?}
\respuesta{Porque más no agrega nada. Con $\eta = 2\pi$ el ruido ya es uniforme
sobre toda la circunferencia: la dirección nueva es completamente aleatoria e
independiente de la de los vecinos. Es el desorden máximo alcanzable, y
aumentar $\eta$ no cambia la distribución resultante.}

\subsection{Sobre la comparación entre modelos}

\pregunta{¿Cuál es exactamente la diferencia entre Vicsek y el votante?}
\respuesta{Una sola: cómo se combinan las orientaciones del vecindario. Vicsek
toma el promedio circular de todos los vecinos; el votante elige uno solo al
azar, con distribución uniforme, y copia su ángulo. Todo lo demás ---el
vecindario, el ruido, la actualización de la posición, las condiciones de
contorno--- es idéntico. Comparten el mismo código salvo esa función.}

\pregunta{¿Qué esperaban que pasara antes de correr las simulaciones?}
\respuesta{Que el votante perdiera el orden a ruidos mucho menores. El
argumento es que promediar $k$ direcciones ruidosas atenúa el ruido por un
factor del orden de $k^{-1/2}$, mientras que copiar a un único vecino lo
transmite íntegro. Los datos lo confirman: a $\rho = 2$, Vicsek cruza
$v_a = 0{,}5$ en $\eta \approx 2{,}9$ y el votante en $\eta \approx 0{,}37$.}

\pregunta{¿De dónde sale el modelo de votante?}
\respuesta{De Loscar, Baglietto y Vázquez, \textit{Physical Review E} 104,
034111 (2021). El modelo de votante clásico es un modelo de opinión donde un
individuo copia la opinión de un vecino al azar; acá la ``opinión'' es la
dirección de movimiento y es multiestado, porque el ángulo es una variable
continua.}

\subsection{Sobre la física de fondo}

\pregunta{¿Esto es una transición de fase de verdad? ¿De qué orden?}
\respuesta{Es una transición orden--desorden controlada por el ruido, con la
polarización como parámetro de orden. Sobre el orden de la transición no nos
pronunciamos: determinarlo requiere un análisis de tamaño finito con varios
valores de $N$ a densidad constante, que no está en el alcance del TP. Nosotros
la caracterizamos, mostramos que existe y ubicamos dónde ocurre; no
clasificamos su orden.}

\pregunta{¿Qué relación tiene esto con el modelo XY?}
\respuesta{La regla de Vicsek es formalmente la de un ferromagneto XY: cada
dirección tiende a alinearse con sus vecinas y el ruido hace de temperatura. La
diferencia es que los espines del XY están fijos en una red y los nuestros se
mueven. Eso es lo relevante: el teorema de Mermin--Wagner prohíbe el orden de
largo alcance en un XY bidimensional a temperatura finita, y el modelo de
Vicsek sí ordena en 2D. El movimiento transporta información de orientación y
genera una interacción de alcance efectivo mucho mayor que $r_c$.}

\pregunta{¿Por qué es interesante que ordene en dos dimensiones?}
\respuesta{Porque en equilibrio no podría. Ese es el aporte del paper de Vicsek
de 1995, y por eso se titula ``un nuevo tipo de transición de fase''. El sistema
está fuera del equilibrio ---es materia activa--- así que Mermin--Wagner no
aplica.}

\subsection{Sobre la presentación misma}

\pregunta{¿Por qué no muestran los valores de los parámetros en la
introducción?}
\respuesta{Porque la guía de la cátedra separa el modelo matemático general
(sección de Introducción) de la descripción del sistema particular que se
simula (sección de Simulaciones). Los valores concretos de $L$, $r_c$, $v$ y las
densidades están en la diapositiva 9.}

\end{document}
